\chapter{Integrasi}

\section{Integrasi Tmux}

\subsection{Konfigurasi dasar}

Tambahkan baris berikut ke berkas \texttt{\textasciitilde/.tmux.conf}:

\begin{lstlisting}[style=toml]
# Tampilkan status PakAI di status bar kanan
set -g status-right "#(pakai tmux)"
set -g status-interval 30

# Opsional: perpanjang lebar status bar
set -g status-right-length 100
\end{lstlisting}

Muat ulang konfigurasi tmux tanpa harus me-restart:

\begin{lstlisting}[style=terminal]
$ tmux source-file ~/.tmux.conf
\end{lstlisting}

\subsection{Konfigurasi dengan tmux-powerline}

\begin{sloppypar}
Jika menggunakan \textit{tmux-powerline} atau tema kustom, sesuaikan
\texttt{status-right} dengan format yang digunakan tema Anda.
\end{sloppypar}

\begin{tipbox}
  Jalankan \cmd{pakai setup tmux} untuk mendapatkan cuplikan konfigurasi tmux
  yang siap tempel (\textit{ready-to-paste}).
\end{tipbox}

\subsection{Persyaratan font}

Ikon penyedia menggunakan karakter \textbf{Nerd Font}. Pastikan terminal dan tmux
menggunakan font Nerd Font seperti:
\begin{itemize}
  \item JetBrains Mono Nerd Font
  \item Fira Code Nerd Font
  \item DejaVu Sans Mono Nerd Font
\end{itemize}

Jika ikon tidak muncul, ikon akan digantikan dengan dua huruf pertama nama
penyedia (misalnya \texttt{cl} untuk Claude).

\section{Integrasi Waybar}

\subsection{Konfigurasi modul}

Tambahkan modul kustom ke berkas konfigurasi Waybar (biasanya
\path{~/.config/waybar/config.jsonc}):

\begin{lstlisting}[style=toml]
"custom/pakai": {
    "exec": "pakai waybar",
    "return-type": "json",
    "interval": 30,
    "tooltip": true,
    "format": "{}",
    "on-click": "pakai dashboard"
}
\end{lstlisting}

\begin{sloppypar}
Tambahkan \texttt{"custom/pakai"} ke array \texttt{modules-right} (atau
\texttt{modules-left}/\texttt{modules-center} sesuai preferensi):
\end{sloppypar}

\begin{lstlisting}[style=toml]
"modules-right": [
    "custom/pakai",
    "clock",
    "battery"
]
\end{lstlisting}

\subsection{Konfigurasi CSS}

Tambahkan style ke berkas CSS Waybar (\texttt{style.css}) untuk
pewarnaan berdasarkan status:

\begin{lstlisting}[style=toml]
#custom-pakai {
    padding: 0 8px;
    color: #cdd6f4;
}

#custom-pakai.ok {
    color: #a6e3a1;
}

#custom-pakai.warning {
    color: #f9e2af;
}

#custom-pakai.critical {
    color: #f38ba8;
}

#custom-pakai.error {
    color: #fab387;
}
\end{lstlisting}

\subsection{Tooltip}

\begin{sloppypar}
Waybar menampilkan tooltip berisi rincian penggunaan per jendela waktu ketika
kursor diarahkan ke modul PakAI. Pastikan \texttt{"tooltip": true} sudah
dikonfigurasi di modul.
\end{sloppypar}

\begin{tipbox}
  Jalankan \cmd{pakai setup waybar} untuk mendapatkan cuplikan konfigurasi
  Waybar yang lengkap, termasuk JSON modul dan CSS.
\end{tipbox}

\section{Integrasi Shell Prompt}

PakAI juga dapat ditampilkan di prompt shell. Contoh untuk Zsh dengan
\textbf{Starship}:

\begin{lstlisting}[style=toml]
# ~/.config/starship.toml
[custom.pakai]
command = "pakai tmux"
when = "pakai daemon status"
format = "[$output]($style) "
style = "green"
\end{lstlisting}
